\section*{Ethics}\label{s:ethical-considerations-appendix}

\emph{Based on the USENIX Security Symposium's policy on ethical considerations}~\cite{secethics}.

Within up to one page, explain the ethical considerations of your work. You are expected to complete a stakeholder-based ethics analysis or justify and complete an alternative \emph{recognized} approach to considering ethics. For a stakeholder-based ethics analysis, you should identify the relevant stakeholders, consider how your work impacts them (positively or negatively), and discuss any steps you have taken to mitigate negative impacts. If your work involves human subjects, you should also discuss how you have ensured that your work complies with relevant ethical guidelines and regulations (e.g., obtaining informed consent, ensuring privacy and confidentiality, etc.). You may also discuss any broader societal implications of your work. If your work does not involve human subjects or have any ethical implications, please provide a brief justification for this conclusion.

For more details on how to complete a stakeholder-based ethics analysis, see, for instance, the USENIX Security Symposium's guidelines on ethical considerations~\cite{secethics}.

\section{Ethical Considerations Appendix}
This section is \textbf{mandatory}. Your ethical considerations go here. Replace this text with your own content.

\section*{AI}\label{s:ethical-considerations-appendix}

\emph{Based on the IEEE Security \& Privacy Symposium's LLM policy}~\cite{spai}.

Students are permitted to use AI (e.g., LLMs) when conducting research and preparing their thesis manuscript. However, they must adhere to the following guidelines to ensure responsible and ethical use of such technologies and detail their usage in the following AI Usage Considerations Appendix.

We ask that students adhere to three key criteria with regards to their use of AI in the scientific process:
\begin{enumerate}
    \item \textbf{Originality}: First, students are responsible for the entire content of their thesis, including all text and figures. While any tool may be used for writing, it is crucial that all content is accurate and original, ensuring transparency and maintaining the integrity of the research process. In particular, students are responsible for the thoroughness of their literature review and must determine relevant prior work and cite it to ensure proper credit. If the students have used AI (e.g., LLMs) to improve their writing, they should, for example, state: ``\emph{LLMs were used for editorial purposes in this manuscript, and all outputs were inspected by the authors to ensure accuracy and originality}.''
    \item \textbf{Transparency}: Second, students should carefully reason about the implications of using AI in their work. If AI is integral to the thesis' methodology, their use should be explicitly detailed. Any idea generated by an AI should be independently developed and validated by the students. Furthermore, students must elaborate on how they handled limitations introduced in their work by their use of AI. Such limitations could for instance include difficulties to obtain results that are reproducible when the AI used is not open sourced.
    \item \textbf{Responsibility}: Third, students should take care to develop LLMs (and AI models in general) responsibly. Any data collection towards training models should take into account relevant ethical considerations such as consent and data holder rights, including intellectual property. Students also have to justify the need for the environmental footprint of their experiments to achieve their goals and support their methodology. We emphasize that the goal here is not to calculate the exact footprint but rather explain experimental choices made as part of the scientific process (e.g., why was an LLM necessary, why was a particular model size selected, how the authors minimized the volume of queries made, which hardware was used to run experiments).
\end{enumerate}

\section{AI Usage Considerations Appendix}
This section is \textbf{mandatory}. Your AI usage considerations go here. Replace this text with your own content.

\section*{Artifacts}\label{s:artifact-appendix}

\emph{Based on the standard Artifact Appendix template used by the USENIX Security
Symposium}.

Please follow the self-contained instructions below and assume
(i)~you are compiling the final version of the Artifact Appendix (no need to
cater to reviewers, but to general users of your artifact) and (ii)~you are
documenting everything in scope for all the artifact evaluation badges
(\emph{Artifacts Available}, \emph{Artifacts Functional}, \emph{Results
Reproduced}). More information at~\cite{secartifacts}.

\input{sections/usenix-ae.tex}

%%% Local Variables:
%%% mode: latex
%%% TeX-master: "../thesis"
%%% End:
